\scanpage{56}

\margin{\begingroup
\newsavebox{\scanFiftySixBubbleBox}
\setbox\scanFiftySixBubbleBox=\hbox{\begin{tabular}{@{}c@{}}
en \unsure{supposant} que les bitors.\\
\unsure{birigidifiés} sur $A,A$ sont des\\
biextensions
\end{tabular}}
\leavevmode
\rlap{\raisebox{-\dp\scanFiftySixBubbleBox}[0pt][0pt]{\begin{tikzpicture}[overlay,baseline=0pt,x={\wd\scanFiftySixBubbleBox},y={\dimexpr\ht\scanFiftySixBubbleBox+\dp\scanFiftySixBubbleBox\relax}]
\draw[line width=.35pt] (.10,1.08) .. controls (.35,1.12) and (.74,1.00) .. (.98,1.08)
 .. controls (1.07,1.06) and (1.07,.82) .. (1.06,.20)
 .. controls (1.03,-.02) and (.89,-.04) .. (.79,-.02)
 .. controls (.74,-.02) and (.75,-.46) .. (.72,-.55)
 .. controls (.68,-.46) and (.70,-.02) .. (.62,.02);
\draw[line width=.35pt] (.59,.02) .. controls (.37,.04) and (.04,.03) .. (-.03,.28)
 .. controls (-.08,.47) and (-.06,.85) .. (-.03,.98);
\draw[line width=.35pt] (.05,1.05) -- (.07,1.06);
\end{tikzpicture}}}
\box\scanFiftySixBubbleBox
\endgroup}
Reprenons la suite \unsure{ex.te} de la fin de \unsure{5|§}), On va définir un hom en sens inverse
\diagram{\begin{tikzcd}[ampersand replacement=\&,column sep=large]
R^1f_*(G_A) \& \underline{\operatorname{Tors}}[\underline{\mathrm{birig}}\ \underline{\mathrm{sym}}](A,A;G)\arrow[l,bend right=15,"\Delta^*"']
\end{tikzcd}}
\margin{\begingroup
\setlength{\unitlength}{0.5pt}
\begin{picture}(470,108)
\put(45,64){inutile pour définir}
\put(285,20){\(\Delta^*\)}
\put(0,108){\begin{tikzpicture}[overlay,x=0.5pt,y=-0.5pt]
\draw (50,57) .. controls (53,74) and (28,92) .. (15,97) .. controls (9,98) and (8,94) .. (11,88);
\draw (5,94)--(11,88)--(17,98);
\end{tikzpicture}}
\end{picture}
\endgroup}
en utilisant l'application diagonale \raisebox{-0.8ex}{\begin{tikzpicture}[x=0.08cm,y=0.08cm,line width=0.5pt]
\draw (0,0) .. controls (5,3) and (8,7) .. (8,9);
\draw (3,-1) .. controls (8,3) and (11,7) .. (10,8) .. controls (8,9) and (5,8) .. (3,8);
\draw (4,7)--(9,12);
\end{tikzpicture}} $\Delta:A\times_S A\mathbin{-}A$.

Calculons le \begingroup
\setlength{\unitlength}{0.5pt}
\begin{picture}(395,167)
\put(0,122){\deleted{\unsure{de}}}
\put(120,104){\inserted{\(h\Delta^*\unsure{j'|g'},\)}}
\put(0,32){composé}
\put(280,32){de}
\put(0,167){\begin{tikzpicture}[overlay,x=0.5pt,y=-0.5pt]
\draw (103,86) .. controls (140,70) and (164,91) .. (174,124) .. controls (198,103) and (220,80) .. (247,61);
\end{tikzpicture}}
\end{picture}
\endgroup $\Delta^*j'$ avec $h:R^1f_*(G_A)\unsure{\mathbin{-}|\longrightarrow}\underline{\operatorname{NS}}(A)$, on aura d'abord
\[
\Delta^*j(L)\simeq(2\operatorname{id}_A)^*(L)L^{-2}
\]
\inserted{\unsure{d'où|Ainsi}} \unsure{Je} dis que c'est
\[
\equiv L^2\quad\{\mathrm{mod}\ A^{\prime}=\underline{\operatorname{Ext}^{\unsure{1}}}(A,G)\}.
\]
Cela résulte du fait suivant

\uline{Lemme} Pour tout entier $n$, $(n\operatorname{id}_A)^*(L)\equiv n^2L$,
\deleted{\unsure{pour $L\in$ des|pour $L\in$ dans}} plus généralement, si $A,B$ sont \unsure{comme|communs} énoncés, alors
\[
\operatorname{Hom}(A,B)\longrightarrow\operatorname{Hom}(\underline{\operatorname{NS}}(B),\underline{\operatorname{NS}}(A))
\]
est une application \uline{quadratique} i.e.
\[
\varphi(u+v)-\varphi(u-v)=2(\varphi(u)+\varphi(v)).
\]
\deleted{\unsure{En ef|Enf}} \uline{Dém} En effet, cette application se factorise en
\diagram{\begin{tikzcd}[ampersand replacement=\&,column sep=small]
\operatorname{Hom}(A,B) \& {\unsure{\mathbin{-}|\longrightarrow}} \& \operatorname{Hom}(\underline{\operatorname{Bitors}}(B,B;G),\underline{\operatorname{Bitors}}(A,A;G))^!\arrow[d,"\text{linéaire}"]\\
 \& \& \operatorname{Hom}(\underline{\operatorname{NS}}(B),\underline{\operatorname{NS}}(A))
\end{tikzcd}}
\unsure{où|à} \unsure{l'exposant|l'argument} ! \unsure{indique|indiquer} les \unsure{hom|liens} qui envoient $\underline{\operatorname{NS}}(B)$ dans $\underline{\operatorname{NS}}(A)$. \deleted{\unsure{On donne|On trouve}} Or il est évident que la première flèche est \unsure{quadratiques|quadratique}..

\margin{55}
